\section{The Divided Will and the Grace to Return}

The need to repent can be clear while the will remains divided. St.~Augustine describes this condition with painful candor in the eighth book of his \emph{Confessions}. By this stage of his account he can no longer shelter behind uncertainty about the truth. He admires the surrender made by others and knows that he ought to give himself to God. Yet he remains attached to the life from which he longs to be freed. His account helps us recognize the interior resistance which can survive even an earnest meditation on eternal punishment.

Augustine traces a chain from disordered willing through indulgence to habit. Repeated indulgence has made the old desire strong, while his newer desire to serve God has not yet overcome it (VIII.5.10). The chain is his own history of consent. It has not appeared from nowhere, and he does not describe himself as merely an innocent spectator of what his desires have become. At the same time, he knows their present force from within. The past has formed a difficulty which a bare announcement of good intentions has failed to remove.

Augustine's account explains why the formation of a habit deserves attention before it has become deeply established. Each consent makes another consent easier, and a repeated excuse becomes readily available when conscience protests. His narrative also gives hope to the person who already knows that difficulty. Augustine writes after his deliverance and remembers how little his own resolutions had accomplished. The same history which exposes his responsibility records what grace has healed. A person recognizing himself in the chain can therefore ask for help without denying either his own past consent or God's power to free him.

He compares his condition to someone struggling to wake from a heavy sleep. He knows that waking is better, but the comfort of remaining asleep keeps defeating the movement to rise. The promise of getting up soon becomes a way of continuing to lie still (VIII.5.12). This image reaches a form of delay more evasive than open denial. We can agree with every warning, intend some future amendment, and continue the same refusal. Agreement then gives us the appearance of conversion while leaving its decisive act indefinitely postponed.

The danger is particularly clear when we think of final impenitence. We need not announce an intention to reject God forever in order to keep rejecting a present call. We can instead reserve today for the attachment and imagine tomorrow as the day in which everything will change. Augustine's recollection lets us hear the weakness of that bargain from inside a life. Time by itself had not supplied the surrender he was withholding. He needed grace to consent now, in the actual struggle, with the desires that made consent difficult.

In the garden with Alypius, the conflict becomes acute. Augustine can command bodily movements and perform them, yet cannot simply command himself into the complete willing he desires. He discovers that part of him wills and part resists (VIII.8.19--9.21). In the next chapter he insists that both movements belong to the same person. He is not two beings, one innocent and one guilty. The divided will is his, and the healing must become his.

His candor can free an examination from vague language. We may say that circumstances prevent a return when we mean that we fear its cost. We may say that we cannot give something up when the truth is that we have not yet consented to the loss. Other impediments may be real and deserve patient help; Augustine is describing his own responsible resistance. The reader should bring the same honesty to his particular case.

Augustine personifies his old pleasures as voices tugging at him while he turns away. They ask whether he can live without them. Their claim is persuasive because they are familiar, while the freedom being offered remains untried (VIII.11.25--26). Temptation thus presents conversion as an empty future deprived of its most necessary goods. The sinner can fear deliverance almost as though it were destruction. A life without the cherished wrong may feel unimaginable precisely because the wrong has occupied the imagination for so long.

Against those voices Augustine sets the figure of Continence, accompanied by the examples of people of different ages and conditions who have received her gift. They have not achieved fidelity by possessing a nature altogether unlike his. God has given them what they could not securely give themselves. The invitation is to stop attempting to stand upon his own strength and entrust himself to the Lord who can receive and heal him (VIII.11.27). The good life appears as fruitful and gracious, and the heart is invited to believe that obedience can be lived.

At the turning point, Augustine hears the repeated summons to take up and read, opens the apostolic book, and receives the command to put on Christ and stop providing for sinful desire. He recounts the inward assurance that followed. Alypius joins him, and Monica receives the news with thanksgiving (VIII.12.28--30). The command addresses the provision Augustine had been making for the desire that divided him. He must cease arranging his life for its continuance. His narrative brings the long argument with himself to a consent which he receives the grace to make.

Augustine has described a struggle in which the danger is known and surrender remains difficult. An annual meditation can expose the same distance between conviction and consent. Its urgency calls the person to stop postponing the duty he acknowledges; Augustine's thanksgiving gives him reason to ask for the grace he needs. The old attachment may still appear indispensable. That judgment has been shaped by repeated indulgence, and Augustine's deliverance shows why it need not determine the future. He who once feared to live without it can remember it afterward as a bondage from which God freed him.

